This study examines when user-controlled devices can share the computation of an AI assistant while retaining explicit data-locality and authorization boundaries. The evaluation compares local execution, an HTTP gateway and cumulative peer-to-peer stage offload under declared tasks, model artifacts, settings and hardware conditions. The current evidence snapshot includes completed typed-request controls and one real prerecorded-audio chain through speech recognition, language-model inference and speech synthesis. The audio chain generated valid output speech but failed the reference task because its recognized instruction differed from the expected input. This negative result is retained separately from the successful typed controls. Full audio-condition measurements and real cross-device comparisons remain incomplete; no audio tail-latency, causal acceleration, deployment-wide privacy or general offload-benefit conclusion is claimed. The contribution at this stage is an implemented sharing architecture and an auditable evaluation method that separates functional correctness, execution locality, resource feasibility and latency. Final conclusions depend on completing the declared non-iOS conditions and reporting their failures and limitations as well as their successes.

\textbf{Keywords:} Distributed Systems; P2P Networks; Data Locality; Voice Assistants.
